%% datasheet.tex — AX101-OD Edge-AI Object Detection Processor
%%
%% Vendored worked example for anvil:datasheet (issue #529). Compile with:
%%   xelatex datasheet.tex   (NOT pdflatex — see anvil-datasheet.cls)
%%
%% Synthesized, NON-CONFIDENTIAL illustrative content — not a real part. Every
%% numeric claim traces to ../refs/spec-bundle.md. The pin-map markers wrap a
%% complete QFN48 pinout (every pin assigned exactly once); the bus-width marker
%% covers the 7-bit ROI index. status: preliminary — every pre-silicon value
%% carries a provenance label and the sheet ends with the preliminary notice.

\documentclass{anvil-datasheet}

%% Signature color override (default 1F4E7A navy lives in the class).
\definecolor{accent}{HTML}{1F4E7A}
\definecolor{rule}{HTML}{1F4E7A}

\datasheetpart{AX101-OD}
\datasheettitle{Edge-AI Object Detection Processor}
\datasheetsubtitle{Single-chip camera-input inference at the edge}
\datasheetcompany{Northgate Silicon}
\datasheetdate{June 2026}
\datasheetrev{0.1}
\datasheetstatus{PRELIMINARY --- SPECIFICATIONS SUBJECT TO CHANGE}

\begin{document}

\maketitleblock

%% ── 1. First page: Key Features | Applications (two-column) ─────────────────
\begin{featurecolumns}
\subsection*{Key Features}
\begin{itemize}
\item Single-shot object detector, 320\texttimes320 RGB input, fully on die
\item \simval{30}\,fps inference at 800\,mV / 600\,MHz (system-model sim.)
\item 2-lane MIPI CSI-2 camera receiver
\item 4\,MB on-die weight SRAM --- no external DRAM required
\item Quad-SPI boot, I\textsuperscript{2}C control, 4\texttimes\ GPIO
\item QFN48, 7\texttimes7\,mm, --40 to +85\,°C
\end{itemize}
\columnbreak
\subsection*{Applications}
\begin{itemize}
\item Smart cameras and video doorbells
\item Retail people-counting and footfall analytics
\item Industrial presence / occupancy detection
\item Battery-powered always-on vision sensors
\end{itemize}
\end{featurecolumns}

\subsection*{General Description}

The AX101-OD is an edge-AI inference processor that runs a quantized single-shot
object detector entirely on die, taking its input directly from a MIPI CSI-2
camera and emitting a list of detected regions of interest with no host
involvement in the inference loop. It is the object-detection SKU of the AX101
family: a single 3.1\,mm² base die in a QFN48 package, configured here for
320\texttimes320 RGB detection. The differentiating fact is the 4\,MB on-die
weight SRAM, which holds the detector's quantized weights with no external DRAM,
making the part a self-contained vision sensor that boots from a small QSPI
flash and runs from a single 800\,mV core supply.

%% ── 2. Device family / Ordering information ─────────────────────────────────
\section{Device Family \& Ordering Information}

\begin{specbox}
\begin{tabularx}{\textwidth}{@{} l X l l l @{}}
\toprule
\textbf{Part number} & \textbf{Function} & \textbf{Package} & \textbf{Temp.\ grade} & \textbf{Status} \\
\midrule
AX101-OD & Object detection (single-shot, 320\texttimes320) & QFN48 & --40 to +85\,°C & Preliminary \\
AX101-OCR & Text recognition (CRNN line OCR) & QFN48 & --40 to +85\,°C & Preliminary \\
\bottomrule
\end{tabularx}
\end{specbox}

The two AX101 SKUs share one base die. \textbf{Shared across the family}
(identical on both sheets): the TSMC 22ULL process, the 3.1\,mm² die, the QFN48
package, the Absolute Maximum Ratings, and the DC / Electrical Characteristics.
\textbf{Per-SKU} (differentiated): the configured network (object detection vs.\
text recognition) and the Performance Characteristics that network delivers. A
customer mixing SKUs in one design can rely on the shared blocks being identical.

%% ── 3. Functional description / block diagram ───────────────────────────────
\section{Functional Description}

The AX101-OD integrates a MIPI CSI-2 receiver, an image signal pre-processor, a
quantized neural inference engine backed by 4\,MB of weight SRAM, and a control
subsystem (QSPI boot, I\textsuperscript{2}C, GPIO) on one die. Camera frames
enter through the 2-lane CSI-2 receiver, are normalized to the 320\texttimes320
detector input, and are streamed into the inference engine; detected regions are
written to an output FIFO the host reads over I\textsuperscript{2}C.

\begin{callout}[title=Block diagram]
A block diagram (CSI-2 RX $\rightarrow$ ISP $\rightarrow$ inference engine + 4\,MB
weight SRAM $\rightarrow$ ROI FIFO; QSPI / I\textsuperscript{2}C / GPIO control)
renders to \texttt{figures/block-diagram.pdf} after \texttt{datasheet-figures}.
This vendored example ships the empty \texttt{figures/} dir; the figure is an
author-supplied artifact stubbed by default.
\end{callout}

The inference engine holds the full detector network in the 4\,MB on-die weight
SRAM (4\,MB; source: model-export), so no external DRAM is required. The CSI-2
receiver is a 2-lane configuration (source: rtl-params).

%% ── 4. Specifications ───────────────────────────────────────────────────────
\section{Specifications}

\subsection{Absolute Maximum Ratings}

\begin{specbox}
\begin{tabularx}{\textwidth}{@{} X l l l l @{}}
\toprule
\textbf{Parameter} & \textbf{Symbol} & \textbf{Min} & \textbf{Max} & \textbf{Unit} \\
\midrule
Supply voltage (core) & $V_{DD}$ & --0.3 & 1.2 & V \\
Supply voltage (I/O) & $V_{DDIO}$ & --0.3 & 3.6 & V \\
Storage temperature & $T_{STG}$ & --55 & 150 & °C \\
\bottomrule
\end{tabularx}
\end{specbox}

\subsection{Recommended Operating Conditions}

\begin{specbox}
\begin{tabularx}{\textwidth}{@{} X l l l l l @{}}
\toprule
\textbf{Parameter} & \textbf{Symbol} & \textbf{Min} & \textbf{Typ} & \textbf{Max} & \textbf{Unit} \\
\midrule
Core supply & $V_{DD}$ & 0.72 & 0.80 & 0.88 & V \\
I/O supply & $V_{DDIO}$ & 1.71 & 1.80 & 1.89 & V \\
Ambient temperature & $T_A$ & --40 & 25 & 85 & °C \\
\bottomrule
\end{tabularx}
\end{specbox}

\subsection{DC / Electrical Characteristics}

\begin{specbox}
\begin{tabularx}{\textwidth}{@{} X l l l l l X @{}}
\toprule
\textbf{Parameter} & \textbf{Symbol} & \textbf{Min} & \textbf{Typ} & \textbf{Max} & \textbf{Unit} & \textbf{Notes} \\
\midrule
Active power & $P_{ACT}$ & --- & \simval{120} & --- & mW & system-model simulation \\
Standby current & $I_{STBY}$ & --- & \est{50} & --- & µA & characterization pending \\
\bottomrule
\end{tabularx}
\end{specbox}

%% ── 5. Performance Characteristics ──────────────────────────────────────────
\clearpage
\section{Performance Characteristics}

Performance is quoted for the shipped single-shot detector network at
320\texttimes320 RGB input, at the nominal 800\,mV / 600\,MHz operating point.
Every value below is pre-silicon and carries its provenance label; the numbers
trace to the system-model simulation in \texttt{refs/spec-bundle.md}.

\begin{specbox}
\begin{tabularx}{\textwidth}{@{} X l l X X @{}}
\toprule
\textbf{Parameter} & \textbf{Typ} & \textbf{Unit} & \textbf{Conditions} & \textbf{Notes} \\
\midrule
Inference rate & \simval{30} & fps & 320\texttimes320 input, 600\,MHz & system-model simulation \\
Single-frame latency & \simval{18} & ms & single frame, cold FIFO & system-model simulation \\
ROI records per frame & 100 & --- & 7-bit \texttt{roi\_index} & rtl-params \\
\bottomrule
\end{tabularx}
\end{specbox}

The detector emits up to 100 region-of-interest records per frame, indexed by a
7-bit field (capacity 128).
% anvil-bus: name=roi_index width=7 max=99

%% ── 6. Pin Configuration and Functions ──────────────────────────────────────
\clearpage
\section{Pin Configuration and Functions}

Pins are numbered 1--48 counter-clockwise from the package top-view fiducial
(pin 1 mark). The QFN48 ballout below assigns every pin exactly once.

\begin{specbox}
\begin{tabularx}{\textwidth}{@{} l l l X @{}}
\toprule
\textbf{Pin} & \textbf{Name} & \textbf{Type} & \textbf{Description} \\
\midrule
% anvil-pinmap-begin package=QFN48 pins=48
1 & VDD\_CORE & P & Core supply (0.80\,V nom.) \\
2 & VSS & G & Ground \\
3 & MIPI\_D0P & I & MIPI CSI-2 data lane 0, positive \\
4 & MIPI\_D0N & I & MIPI CSI-2 data lane 0, negative \\
5 & MIPI\_D1P & I & MIPI CSI-2 data lane 1, positive \\
6 & MIPI\_D1N & I & MIPI CSI-2 data lane 1, negative \\
7 & MIPI\_CKP & I & MIPI CSI-2 clock lane, positive \\
8 & MIPI\_CKN & I & MIPI CSI-2 clock lane, negative \\
9 & VDDIO & P & I/O supply (1.80\,V nom.) \\
10 & VSS & G & Ground \\
11 & QSPI\_CLK & O & Quad-SPI boot clock \\
12 & QSPI\_CS & O & Quad-SPI chip select \\
13 & QSPI\_IO0 & I/O & Quad-SPI data 0 \\
14 & QSPI\_IO1 & I/O & Quad-SPI data 1 \\
15 & QSPI\_IO2 & I/O & Quad-SPI data 2 \\
16 & QSPI\_IO3 & I/O & Quad-SPI data 3 \\
17 & VDD\_CORE & P & Core supply (0.80\,V nom.) \\
18 & VSS & G & Ground \\
19 & I2C\_SCL & I/O & I\textsuperscript{2}C control clock \\
20 & I2C\_SDA & I/O & I\textsuperscript{2}C control data \\
21 & GPIO0 & I/O & General-purpose I/O 0 \\
22 & GPIO1 & I/O & General-purpose I/O 1 \\
23 & GPIO2 & I/O & General-purpose I/O 2 \\
24 & GPIO3 & I/O & General-purpose I/O 3 \\
25 & VDDIO & P & I/O supply (1.80\,V nom.) \\
26 & VSS & G & Ground \\
27 & INT\_N & O & Frame-ready interrupt (active low) \\
28 & RESET\_N & I & Chip reset (active low) \\
29 & BOOT\_SEL & I & Boot source strap \\
30 & CLK\_REF & I & Reference clock input \\
31 & VDD\_CORE & P & Core supply (0.80\,V nom.) \\
32 & VSS & G & Ground \\
33 & VDD\_PLL & P & PLL analog supply \\
34 & VSS\_PLL & G & PLL analog ground \\
35 & TEST\_EN & I & Test-mode enable (tie low in application) \\
36 & TDI & I & JTAG test data in \\
37 & TDO & O & JTAG test data out \\
38 & TCK & I & JTAG test clock \\
39 & TMS & I & JTAG test mode select \\
40 & TRST\_N & I & JTAG test reset (active low) \\
41 & VDDIO & P & I/O supply (1.80\,V nom.) \\
42 & VSS & G & Ground \\
43 & VDD\_CORE & P & Core supply (0.80\,V nom.) \\
44 & VSS & G & Ground \\
45 & VDD\_SRAM & P & Weight-SRAM supply \\
46 & VSS\_SRAM & G & Weight-SRAM ground \\
47 & SPARE0 & --- & Reserved --- leave unconnected \\
48 & SPARE1 & --- & Reserved --- leave unconnected \\
% anvil-pinmap-end
\bottomrule
\end{tabularx}
\end{specbox}

%% ── 7. Typical Application ──────────────────────────────────────────────────
\section{Typical Application}

A reference integration drives a MIPI CSI-2 camera module into the 2-lane
receiver, boots the AX101-OD from a small QSPI flash holding the detector
weights, and reads ROI records over I\textsuperscript{2}C. The part runs from a
single 800\,mV core supply and a 1.80\,V I/O supply; the weight SRAM supply
(\texttt{VDD\_SRAM}) is decoupled separately. Hold \texttt{TEST\_EN} low and
strap \texttt{BOOT\_SEL} for QSPI boot in the application.

\begin{callout}[title=Reset and boot]
On power-up, hold \texttt{RESET\_N} low until both supplies are stable, then
release; the part fetches its weights from QSPI flash and asserts
\texttt{INT\_N} on the first detected frame. The 320\texttimes320 input size
quoted in Performance Characteristics is the size the host's camera pipeline
must deliver --- it agrees with the inference-rate conditions above.
\end{callout}

%% ── 8. Package / Mechanical ─────────────────────────────────────────────────
\section{Package \& Mechanical}

QFN48, 7\texttimes7\,mm body, 0.5\,mm pitch, 0.40\,mm nominal terminal width,
with an exposed center pad tied to \texttt{VSS}. The 48-pin ballout is consistent
with the family / ordering table and the pinout above (the auditor cross-checks
the three). The package outline traces to the OSAT drawing referenced in
\texttt{refs/}.

\begin{callout}[title=Package outline]
A dimensioned package-outline drawing renders to
\texttt{figures/package-outline.pdf}; it is an author-supplied artifact stubbed
by default in this vendored example.
\end{callout}

%% ── 9. Revision History ─────────────────────────────────────────────────────
\section{Revision History}

\begin{specbox}
\begin{tabularx}{\textwidth}{@{} l l X @{}}
\toprule
\textbf{Rev} & \textbf{Date} & \textbf{Changes} \\
\midrule
0.1 & June 2026 & Initial preliminary release. \\
\bottomrule
\end{tabularx}
\end{specbox}

%% ── 10. Legal / notices ─────────────────────────────────────────────────────
\preliminarynotice

{\footnotesize\color{muted} Information in this document is provided in
connection with the named products. No license, express or implied, to any
intellectual property rights is granted by this document. The vendor reserves
the right to make changes to specifications and product descriptions at any
time. This is a synthesized, NON-CONFIDENTIAL worked example --- not a real
part.}

\end{document}
